\chapter{System Architecture}
\label{ch:architecture}

\section{Design Principles}
\label{sec:arch-principles}

Three principles shape the implementation described in the remaining
chapters.

\begin{enumerate}
    \item \textbf{Immutability of the circuit model.} A \code{Circuit}
          object, once constructed from a netlist and a specific assignment
          of cell sizes, is never mutated. Its topological order, per-gate
          load capacitance, per-gate delay, total area, and total power are
          computed once, eagerly, at construction time. Resizing a gate does
          not modify an existing \code{Circuit}; it produces an entirely new
          \code{Circuit} instance with one gate's cell replaced and every
          net reference correctly rebuilt. This trades additional memory
          allocation during optimization (\cref{ch:optimization}, where many
          candidate resizes are evaluated per iteration) for the elimination
          of an entire class of bugs in which a trial mutation is
          incompletely reverted and silently corrupts subsequent analysis.
    \item \textbf{Separation of static circuit state from one analysis
          run's results.} A \code{Circuit} describes \emph{what the circuit
          currently is}: its topology and its currently-assigned sizes. It
          does not store arrival times, required times, or slack. Those are
          the output of a separate, pure \code{analyze\_timing} function
          that takes a \code{Circuit} and a configuration as input and
          returns a timing-result value. This means the same
          \code{Circuit} can be analyzed multiple times (for example, once
          per Monte Carlo sample, or once per stage in a diagnostic sweep)
          without the results of one analysis leaking into another.
    \item \textbf{Pure, side-effect-free analysis modules.} STA, logical
          effort, power/area estimation, and Monte Carlo analysis are each
          implemented as functions that take immutable inputs and return a
          new result value, rather than as methods that accumulate state on
          an object across calls. This makes each module independently
          testable and makes the overall data flow
          (\cref{sec:arch-dataflow}) easy to reason about: every stage's
          inputs are exactly the outputs of the stages before it.
\end{enumerate}

\section{Module Map}
\label{sec:arch-modules}

The implementation is organized into eight top-level packages, listed here
in the order a reader unfamiliar with the codebase should study them ---
each package depends only on the ones listed above it.

\begin{longtable}{@{}p{0.24\textwidth}p{0.70\textwidth}@{}}
\toprule
\textbf{Package} & \textbf{Responsibility} \\
\midrule
\endhead
\code{domain/} &
  Owns the \code{Cell}, \code{Gate}, \code{Net}, and \code{Circuit} data
  model, and shared numeric helpers. Every other package depends on this
  one; it depends on nothing else in the project. \\[4pt]
\code{input/} &
  Parses and validates the netlist, standard-cell library, and
  configuration file, and builds the initial \code{Circuit}
  (\cref{ch:validation}). \\[4pt]
\code{analysis/} &
  Static timing analysis, logical-effort path
  analysis, and Monte Carlo statistical analysis
  (\cref{ch:sta,ch:logical-effort,ch:monte-carlo}). \\[4pt]
\code{optimization/} &
  The cost function, the four gate-sizing heuristics, and the optimizer
  driver loop (\cref{ch:optimization}). \\[4pt]
\code{reporting/} &
  Serializes every analysis result to the on-disk output contract (CSV,
  JSON, and PNG artifacts), including the topology export consumed by the
  interactive viewer. \\[4pt]
\code{app/} &
  Coordinates a full single-circuit run end to end and exposes the
  unified command-line interface. \\[4pt]
\code{benchmarking/} &
  The synthetic benchmark generation and evaluation framework
  (\cref{ch:benchmarking}). \\[4pt]
\code{viewer/} &
  The interactive topology viewer's local server and static front-end
  assets (\cref{ch:viewer}). \\
\bottomrule
\end{longtable}

\section{Key Classes and Public API}
\label{sec:arch-api}

Before tracing the data flow (\cref{sec:arch-dataflow}), it is worth naming
the concrete classes and functions each phase is built from, since the
remainder of this report refers to them by name rather than only
describing their behavior in prose. \Cref{tab:arch-domain-api} lists the
core data model (\code{domain/}); the analysis, optimization, and reporting
entry points are listed at the start of their respective chapters
(\cref{ch:sta,ch:logical-effort,ch:optimization,ch:monte-carlo}).

\begin{longtable}{@{}p{0.30\textwidth}p{0.62\textwidth}@{}}
\toprule
\textbf{Class / function} & \textbf{Role} \\
\midrule
\endhead
\code{Cell}, \code{InputPin}, \code{RiseFall} &
  Frozen dataclasses holding one library cell's characterization data:
  \code{name}, \code{family}, \code{size}, \code{size\_factor},
  \code{num\_inputs}, \code{input\_pins}, \code{timing\_sense},
  \code{intrinsic\_delay}, \code{output\_resistance},
  \code{internal\_capacitance}, \code{parasitic\_delay},
  \code{leakage\_power}, \code{area}. Constructed and validated by
  \code{Cell.from\_dict(name, data)}. \\[4pt]
\code{CellLibrary(Mapping[\newline str, Cell])} &
  Parses \code{cell\_library.json} in \code{\_\_init\_\_(path)} and exposes
  cells as a read-only mapping. Key methods: \code{find(name)},
  \code{variants(family)}, and \code{sizing\_candidates(family,\allowbreak
  input\_pin,target\_ capacitance)} (\cref{sec:le-candidates}). \\[4pt]
\code{Net}, \code{Gate} &
  Frozen dataclasses describing structural connectivity, independent of
  which cell size is assigned: \code{Net} holds \code{driver} and
  \code{loads: tuple[(pin, gate\_name), ...]}; \code{Gate} holds
  \code{cell\_family}, \code{inputs}, \code{output}. \\[4pt]
\code{PathStep}, \code{CircuitPath} &
  \code{PathStep(gate, input\_pin)} is one gate traversal through a
  specific pin; \code{CircuitPath} is an ordered tuple of steps between a
  primary input and a primary output, exposing \code{.gates} as a derived
  property. \\[4pt]
\code{NetListParser} &
  \code{\_\_init\_\_(netlist\_path, cell\_library)}; \code{parse()} returns
  \code{(nets, gates, cell\_assignment)} after performing every structural
  check of \cref{tab:error-classes}. \\[4pt]
\code{CircuitTopology} &
  Frozen, cell-assignment-independent structure
  (\code{netlist}, \code{gates}, \code{topological\_order}), built once by
  \code{CircuitTopology.create(netlist, gates)} and shared by every sizing
  variant of the same circuit. \\[4pt]
\code{Circuit} &
  One immutable cell-sizing assignment over a shared topology.
  \code{\_\_init\_\_(netlist, gates, gate\_cells, config, cell\_library)}
  eagerly computes area, power, fan-out capacitances, and gate delays.
  Query methods: \code{cell\_for(gate)}, \code{input\_net(gate,pin)},
  \code{output\_net(gate)}. \code{with\_gate\_cell(gate\_name, replacement)}
  returns a new \code{Circuit} with one gate resized, reusing the shared
  topology and updating only the affected metrics incrementally
  (\cref{sec:opt-eval}). \\[4pt]
\code{replace\_gate\_cell(\newline circuit, gate\_name, replacement)} &
  Module-level convenience wrapper around
  \code{Circuit.with\_gate\_cell}. \\[4pt]
\code{replacement\_area\_and \newline \_power(circuit, gate\_name, replacement)} &
  Returns the \emph{delta}-updated $(A_{\mathrm{total}}, P_{\mathrm{total}})$
  a candidate resize would produce \emph{without} constructing a full
  \code{Circuit} variant --- used by the optimizer as a cheap pre-filter before paying for a
  full STA re-run. \\[4pt]
\code{Config} &
  \code{\_\_init\_\_(path)} validates \code{config.json} in full
  (\cref{sec:val-config}) and exposes typed sub-configurations
  (\code{TimingAnalysisConfig}, \code{OperatingConditions},
  \code{DesignConstraints}, \code{OptimizationConfig},
  \code{MonteCarloConfig}) plus per-net accessor methods
  \code{output\_load(net)}, \code{input\_arrival(net)},
  \code{output\_required(net)}, \code{activity\_factor(node)}. \\
\bottomrule
\caption{Core data-model classes in \texttt{domain/} and \texttt{input/},
used throughout the rest of this report.}
\label{tab:arch-domain-api}
\end{longtable}

\section{End-to-End Data Flow}
\label{sec:arch-dataflow}

\Cref{fig:dataflow} shows the complete pipeline for a single circuit, from
its two input files through to the reported outputs, annotated with the
governing equation at each stage. The pipeline has five clearly separated
phases: (1) parsing and structural validation (\code{NetListParser.parse},
\code{Config.\_\_init\_\_}), which either rejects the
circuit outright or produces a validated \code{Circuit}; (2) static timing
analysis (\code{analyze\_timing}, \cref{ch:sta}), itself split into a forward pass (arrival time), a backward pass
(required time), and their combination into slack, WNS/TNS, and the top-$K$
critical paths; (3) logical-effort analysis
(\code{analyze\_path\_logical\_effort}, \cref{ch:logical-effort}) of those critical paths; (4) the
optimization loop (\code{CircuitOptimizer.optimize}), which repeatedly proposes a gate resize, evaluates it
with a full re-run of phase (2), and accepts or rejects it according to the
cost function and termination rules of \cref{ch:optimization}; and (5)
Monte Carlo statistical analysis (\code{run\_monte\_carlo}) and report generation. Phases (1)--(5) are
driven end to end by \code{Experiments.run(circuit)}
(\code{app/experiments.py}), which dispatches to one handler function per
named experiment in the fixed order given by \code{DEFAULT\_EXPERIMENTS}:
\code{fanout\_capacitances}, \code{gate\_delays}, \code{timing\_analysis},
\code{logical\_ effort\_analysis}, \code{optimization},
\code{topology\_export}, \code{visualization}, \code{monte\_carlo}, and
\code{summary} --- an explicit, literal ordering that doubles as the most
authoritative statement of the pipeline's dependency order.

\begin{figure}[H]
    \centering
    \includegraphics[width=0.97\textwidth]{Architecture}
    \caption{End-to-end data flow of the analyzer, from input files to
    reported outputs.}
    \label{fig:dataflow}
\end{figure}

\section{Command-Line Interface and Output Contract}
\label{sec:arch-cli}

The command-line entry point is \code{app/cli.py:main(argv)}, which builds
its parser via \code{build\_argument\_parser()} and delegates to
\code{app/application.py:run\_application(argv)}. That function performs,
in order: \code{setup\_logging(verbose, log\_path)}; input loading via
\code{load\_circuit(arguments)} (which internally parses the configuration,
cell library, and netlist, and constructs the initial \code{Circuit}); and,
if loading succeeds, \code{run\_experiments(circuit, netlist\_path,
output\_dir)}, which constructs an \code{Experiments} instance
(\cref{sec:arch-dataflow}) and calls its \code{run(circuit)} method. A
single circuit is analyzed with:

\begin{lstlisting}[language=bash, style=pythonstyle, numbers=none]
./run_circuit.sh c01_inverter_chain
\end{lstlisting}

which writes an independent, self-contained directory
\code{outputs/c01\_inverter\_chain/} containing every artifact listed in
\cref{tab:output-contract}. Invalid inputs produce only
\code{validation\_report .json} and \code{run.log}, and the process exits
with a nonzero status without attempting any timing analysis, so that a
downstream script (or the batch mode, \code{run\_circuit.sh --all}) can
mechanically distinguish an analyzed circuit from a rejected one.

\begin{longtable}{@{}p{0.34\textwidth}p{0.60\textwidth}@{}}
\toprule
\textbf{Artifact} & \textbf{Contents} \\
\midrule
\endhead
\code{validation\_report.json}, \code{run.log} &
  Structured input-validation result and the complete execution log. \\[3pt]
\code{fanout\_capacitances.csv}, \code{gate\_delays.csv} &
  One row per gate, topologically ordered: output load and rise/fall
  delay. \\[3pt]
\code{timing\_analysis.csv} &
  One row per gate: load, rise/fall delay, arrival time, required time,
  transition slack, and overall node slack. \\[3pt]
\code{timing\_summary.csv}, \code{critical\_paths.csv} &
  Circuit delay, WNS, TNS, compliance flag, and the top-$K$ critical
  paths. \\[3pt]
\code{logical\_effort\_paths/ stages/candidates.csv} &
  Path-level ($G,B,H,F,P$), stage-level, and discrete-candidate logical
  effort records. \\[3pt]
\code{optimization\_*.csv}, \code{optimization\_comparison.csv} &
  Complete iteration histories for every heuristic and their head-to-head
  comparison. \\[3pt]
\code{circuit\_graph\_\{pre,post\}\_ optimization.png} &
  Slack-colored circuit diagrams with the critical path highlighted. \\[3pt]
\code{circuit\_topology.json} &
  Versioned graph contract consumed by the interactive viewer
  (\cref{ch:viewer}). \\[3pt]
\code{monte\_carlo\_*.csv}, \code{monte\_carlo\_summary.json} &
  Paired pre-/post-optimization statistical timing results
  (\cref{ch:monte-carlo}), or an explicit ``skipped'' summary when
  disabled. \\[3pt]
\code{summary.json} &
  Aggregate summary for the circuit. \\
\bottomrule
\caption{Output contract produced for every successfully analyzed
circuit.}
\label{tab:output-contract}
\end{longtable}

\newpage
All CSV column names carry explicit units (\code{\_ns}, \code{\_fF},
\code{\_uW}); all numeric computation and serialization retain full
floating-point precision, with rounding applied only when a value is
formatted for a plot label, so that downstream tooling never operates on
already-rounded intermediate values.

